\chapter{堆疊與物件}

堆疊（Stack）是最重要的資料結構之一。 LIFO（後進先出）原則在括號匹配、函式呼叫、表達式求值等場景中無所不在。 這章從堆疊的基本概念開始，介紹兩種實現方式（陣列和鏈結序列），然後展示如何用堆疊解決實際問題。 最後強調用 struct 和函式封裝資料的物件導向優勢。

\section{堆疊概念與操作}

堆疊是後進先出（LIFO）的資料結構。 只能從頂端插入和刪除元素。

\subsection*{堆疊的定義}

\begin{code}{c}{stack basics}
// push: insert at top
// pop: remove from top
// peek: view top without removing
// isEmpty: check if empty
// isFull: check if full

Stack: [1, 2, 3, 4, 5]
              ^
              top
\end{code}

\subsection*{堆疊的應用}

\begin{itemize}
\item 函式呼叫堆疊：自動管理函式的本地變數
\item 括號匹配：檢查括號是否平衡
\item 撤銷功能：存儲編輯歷史
\item 表達式求值：計算 1 + 2 * 3
\item 深度優先搜尋（DFS）：手工模擬遞迴呼叫
\end{itemize}

\keypoint{LIFO 特性使堆疊適合追蹤狀態與回溯問題。}

\section{堆疊的陣列實作}

\subsection*{堆疊結構與標頭檔}

用陣列實現堆疊。 用 \texttt{top} 紀錄頂端位置。

\begin{code}{c}{stack structure}
#define MAX_SIZE 100

typedef struct {
    int data[MAX_SIZE];
    int top;
} Stack;

// stack.h: function declarations
void stack_init(Stack *s);
void push(Stack *s, int value);
int pop(Stack *s);
int peek(Stack *s);
int isEmpty(Stack *s);
int isFull(Stack *s);
\end{code}

\subsection*{堆疊函式實作}

\begin{code}{c}{stack operations}
void stack_init(Stack *s) {
    s->top = -1;
}

void push(Stack *s, int value) {
    assert(s->top < MAX_SIZE - 1);
    s->data[++s->top] = value;
}

int pop(Stack *s) {
    assert(s->top >= 0);
    return s->data[s->top--];
}

int peek(Stack *s) {
    assert(s->top >= 0);
    return s->data[s->top];
}

int isEmpty(Stack *s) {
    return s->top < 0;
}

int isFull(Stack *s) {
    return s->top >= MAX_SIZE - 1;
}
\end{code}

\subsection*{分別編譯方法}

\begin{code}{c}{separate compilation}
// Compile stack.c
gcc -c stack.c -o stack.o

// Compile main.c
gcc -c main.c -o main.o

// Link object files
gcc stack.o main.o -o stack_demo

// Or one-step compilation
gcc stack.c main.c -o stack_demo
\end{code}

\keypoint{分別編譯提高開發效率，支援大型專案。 標頭檔隱藏實現細節。}

\section{括號平衡判別}

\subsection*{問題描述}

判定一個字符串中的括號是否平衡。

\begin{code}{c}{parenthesis examples}
"(())"      // balanced
"(()(("     // not balanced
")("        // not balanced
"()()"      // balanced
\end{code}

\subsection*{演算法}

遇到左括號入堆疊，遇到右括號檢查堆疊頂。

\begin{example}{c}{(balance.c) Parenthesis matching}
#include <stdio.h>
#include <string.h>
#include <assert.h>

#define MAX_SIZE 100

typedef struct {
    char data[MAX_SIZE];
    int top;
} Stack;

void stack_init(Stack *s) {
    s->top = -1;
}

void push(Stack *s, char c) {
    assert(s->top < MAX_SIZE - 1);
    s->data[++s->top] = c;
}

char pop(Stack *s) {
    assert(s->top >= 0);
    return s->data[s->top--];
}

int isEmpty(Stack *s) {
    return s->top < 0;
}

int isBalanced(const char *str) {
    Stack s;
    stack_init(&s);

    for (int i = 0; str[i]; i++) {
        if (str[i] == '(') {
            push(&s, '(');
        } else if (str[i] == ')') {
            if (isEmpty(&s)) return 0;
            pop(&s);
        }
    }
    return isEmpty(&s);
}

int main() {
    printf("%d\n", isBalanced("(())"));
    printf("%d\n", isBalanced("(()"));
    printf("%d\n", isBalanced("()()"));
}
\end{example}

\iosample{
}{
1 \\
0 \\
1 \\
}

\subsection*{多種括號}

擴展支援多種括號類型。

\begin{code}{c}{multiple bracket types}
int isBalancedMulti(const char *str) {
    Stack s;
    stack_init(&s);

    for (int i = 0; str[i]; i++) {
        char c = str[i];
        if (c == '(' || c == '[' || c == '{') {
            push(&s, c);
        } else if (c == ')' || c == ']' || c == '}') {
            if (isEmpty(&s)) return 0;
            char top = pop(&s);
            if ((c == ')' && top != '(') ||
                (c == ']' && top != '[') ||
                (c == '}' && top != '{')) {
                return 0;
            }
        }
    }
    return isEmpty(&s);
}
\end{code}

\keypoint{堆疊自然適合括號匹配問題。}

\section{對稱字符串}

\subsection*{Palindrome 判別}

字符串正讀和倒讀相同。

\begin{example}{c}{(palindrome.c) Check palindrome}
#include <stdio.h>
#include <string.h>
#include <assert.h>

#define MAX_SIZE 100

typedef struct {
    char data[MAX_SIZE];
    int top;
} Stack;

void stack_init(Stack *s) {
    s->top = -1;
}

void push(Stack *s, char c) {
    assert(s->top < MAX_SIZE - 1);
    s->data[++s->top] = c;
}

char pop(Stack *s) {
    assert(s->top >= 0);
    return s->data[s->top--];
}

int isEmpty(Stack *s) {
    return s->top < 0;
}

int isPalindrome(const char *str) {
    int len = strlen(str);
    Stack s;
    stack_init(&s);

    for (int i = 0; i < len / 2; i++) {
        push(&s, str[i]);
    }

    int start = (len % 2 == 0) ? len / 2 : len / 2 + 1;

    for (int i = start; i < len; i++) {
        if (pop(&s) != str[i]) return 0;
    }
    return isEmpty(&s);
}

int main() {
    printf("%d\n", isPalindrome("racecar"));
    printf("%d\n", isPalindrome("madam"));
    printf("%d\n", isPalindrome("hello"));
}
\end{example}

\iosample{
}{
1 \\
1 \\
0 \\
}

\section{堆疊的鏈結序列實作}

\subsection*{動態堆疊}

用鏈結序列實現堆疊，無大小限制。

\begin{code}{c}{stack with linked list}
typedef struct Node {
    int data;
    struct Node *next;
} Node;

typedef struct {
    Node *top;
} Stack;

void stack_init(Stack *s) {
    s->top = NULL;
}

void push(Stack *s, int value) {
    Node *new_node = (Node *)malloc(sizeof(Node));
    assert(new_node != NULL);
    new_node->data = value;
    new_node->next = s->top;
    s->top = new_node;
}

int pop(Stack *s) {
    assert(s->top != NULL);
    Node *temp = s->top;
    int value = temp->data;
    s->top = s->top->next;
    free(temp);
    return value;
}

int isEmpty(Stack *s) {
    return s->top == NULL;
}
\end{code}

\subsection*{介面統一}

陣列版和鏈表版的標頭檔相同。

\begin{code}{c}{same interface}
// Both implementations have same interface:
Stack *stack_create();
void push(Stack *s, int value);
int pop(Stack *s);
int isEmpty(Stack *s);
void stack_destroy(Stack *s);
\end{code}

\keypoint{介面統一是多態的基礎。}

\section{物件導向的優勢}

\subsection*{堆疊作為物件}

用結構和函式封裝資料，形成完整的「物件」。

\begin{code}{c}{stack as object}
// stack.h defines the interface
typedef struct Stack Stack;

Stack *stack_create(int max_size);
void push(Stack *s, int value);
int pop(Stack *s);
void stack_destroy(Stack *s);

// main.c uses the object
Stack *s = stack_create(100);
push(s, 10);
push(s, 20);
printf("%d\n", pop(s));
stack_destroy(s);
\end{code}

\subsection*{OOP 優勢}

\begin{itemize}
\item 模組化：堆疊獨立於主程式
\item 可重用性：同一堆疊可用於多個專案
\item 易於維護：修改實作不影響使用者
\item 易於測試：獨立測試堆疊功能
\item 易於擴展：新增操作只需修改實作檔
\end{itemize}

\keypoint{即使 C 也能實現物件的核心思想。}
